% !PRINTED 84
\weberpage{84}
If this condition is satisfied, then $g_1, g_2, \dots g_\nu$ is certainly the basis of a group with the invariants $i_1, i_2, \dots i_\nu$. But if the requirement set up in IV. is also to be satisfied, then it follows further:

\gesperrt{3. If $q_h$ is a prime dividing $m$, to which the index $\beta_h$ corresponds, then the $\nu$ quantities}
\[
\gamma_{1,h}, \gamma_{2,h}, \dots \gamma_{\nu,h},
\]
\gesperrt{must not, if $q_h$ divides $m$ more than once, all be divisible by $q_h$, and, if $q_h$ divides $m$ only once, must not all be divisible by $q_h-1$.}

The conditions 2., 3. are then accordingly the necessary and sufficient ones, and it now only remains to bring them into a form in which the possibility of their fulfilment can be judged.

We denote by $\delta_{k,h}$ the \gesperrt{greatest common divisor of $i_k$ and $c_h$ and set}
\weq{11}{i_k=\delta_{k,h}\,i_{k,h}, \quad c_h=\delta_{k,h}\,c_{k,h}, \quad \begin{aligned}k&=1,2,\dots\nu,\\ h&=1,2,\dots\mu,\end{aligned}}
so that $i_{k,h}$ and $c_{k,h}$ are \gesperrt{relatively prime}. It is then to be remarked that the quantities
\[
\delta_{k,h}, \quad i_{k,h}, \quad c_{k,h}
\]
are completely determined by $m$ and the invariants $i_k$. By the definition of the invariants, the $\delta_{k,h}$ and $i_{k,h}$ must be \gesperrt{prime powers}.

From the assumptions which we made in a) to d) about the number $m$, there results a consequence for the $c_{k,h}$ which we must emphasise.

\gesperrt{4. If $q_h^\varkappa$ is one of the factors of $m$ to which the index $\beta_h$ corresponds, so that $c_h=q_h^{\varkappa-1}(q_h-1)$, or, if $q=2$ and $\varkappa>2$, $c_h=2^{\varkappa-2}$, and if $\varkappa=2$, $c_h=2$, then, if $\varkappa>1$, the numbers}
\weq{12}{c_{1,h}, c_{2,h}, \dots c_{\nu,h}}
\gesperrt{cannot all be divisible by $q_h$, and, if $\varkappa=1$, cannot all be divisible by $q-1$.}
% [sic: Weber switches from q_h to q in the exceptional 2-adic clause and in the final q-1 clause; q_h is intended in both places.]

For if the quantities (12) are all divisible by $q_h$, then by (11) the numbers
\weq{13}{\delta_{1,h}.\ \delta_{2,h}, \dots \delta_{\nu,h}}
% [sic: Weber prints a full stop, rather than a comma, after the first term.]
%CONT>>
